% from main.tex
%%%%%%%%%%  scan idx 126  |  folio 113  %%%%%%%%%%
une
suite exacte de \(\mathcal C_S\). Si \(G\) est de type fini et \(H\) précohérent, \(F\) est de type fini.

\textbf{b)} Soit
\[
  F\longrightarrow G\longrightarrow H\longrightarrow 0
\]
une suite exacte de \(\mathcal C_S\). Si \(F\) est de type fini et \(G\) précohérent, \(H\) est précohérent.

\textbf{c)} Soit
\[
  0\longrightarrow F\longrightarrow G\longrightarrow H\longrightarrow 0
\]
une suite exacte de \(\mathcal C_S\). Si \(F\) et \(H\) sont précohérents (resp. de type fini), \(G\) l'est aussi.

\emph{Preuve.} a) Quitte à localiser, on a un épimorphisme \(G'\to G\) avec \(G'\in\Ob\mathcal C_0\). D'où un diagramme commutatif dont les lignes sont exactes
\[
\begin{tikzcd}
0\arrow[r]&F'\arrow[r]\arrow[d]&G'\arrow[r]\arrow[d]&H\arrow[d,"\Id"]\\
0\arrow[r]&F\arrow[r]&G\arrow[r]&H .
\end{tikzcd}
\]
\(F'\) est de type fini puisque \(H\) est précohérent, mais \(F'\to F\) est un épimorphisme, donc \(F\) est de type fini.

b) Soient \(U\) au-dessus de \(S\), et \(H'\xrightarrow{f}H|U\) un morphisme de \(\mathcal C_U\) avec \(H'\in\Ob\mathcal C_{0,U}\). Quitte à changer les notations, on peut supposer \(U=S\). Grâce à la quasi-relevabilité de \(\mathcal C_0\) dans \(\mathcal C\), il existe (quitte à localiser) un carré commutatif
\[
\begin{tikzcd}
G' \arrow[r,"p"] \arrow[d] & H' \arrow[d,"f"]\\
G \arrow[r] & H,
\end{tikzcd}
\]
